% !TEX root = 02.tex
\documentclass[a4paper,10pt]{article}

\usepackage{pdfsync}
\usepackage[utf8]{inputenc}
\usepackage{microtype}
\usepackage{mathtools}
\usepackage{amsthm,amsfonts,amsmath,amssymb}
\usepackage{mathabx}
% \usepackage{MnSymbol}
\usepackage{hyperref}
\usepackage{bbold}
\usepackage[textwidth=2cm,textsize=small]{todonotes}
\usepackage{subcaption}
\usepackage{enumerate}
\usepackage{xspace}
\usepackage{stmaryrd}
\usepackage{verbatim}
\usepackage{cleveref}

\hypersetup{
    colorlinks,
    linkcolor={red!50!black},
    citecolor={blue!50!black},
    urlcolor={blue!80!black}
}

\author{}

\newcommand{\ignore}[1]{}

\newcommand{\st}{s.t.}
\newcommand{\ie}{i.e.}
\newcommand{\eg}{e.g.}

\newcommand{\sep}{\;|\;}
\renewcommand{\rule}[2]{\dfrac{#1}{#2}}
\newcommand{\set}[1]{\left\{#1\right\}}
\newcommand{\setof}[2]{\set{#1 \;\middle|\; #2}}
\newcommand{\tuple}[1]{\left( #1 \right)}
\newcommand{\pair}[2]{\tuple {#1, #2}}
\newcommand{\closure}[3]{\tuple {#1, #2, #3}}
\newcommand{\sem}[1]{\llbracket{#1}\rrbracket}
\newcommand{\powerset}[1]{\mathcal P(#1)}
\newcommand{\map}[2]{\mathsf{map}\; #1\; #2}
\newcommand{\filter}[2]{\mathsf{filter}\; #1\; #2}
\newcommand{\floor}[1]{\left\lfloor #1 \right\rfloor}
\newcommand{\PreRel}[2]{\mathsf{Pre}_{#1}(#2)}
\newcommand{\Pre}[1]{\PreRel {} {#1}}
\newcommand{\PreStar}[1]{\mathsf{Pre}^*(#1)}
\newcommand{\Post}[1]{\mathsf{Post}(#1)}
\newcommand{\PostStar}[1]{\mathsf{Post^*}(#1)}
\newcommand{\Conf}{\mathsf{Conf}}
\newcommand{\upwardclosure}[1]{#1\uparrow}
\newcommand{\card}[1]{\left| #1 \right|}
\newcommand{\lang}[1]{L(#1)}
\newcommand{\goesto}[1]{\xrightarrow{#1}}
\newcommand{\e}{\varepsilon}

\newcommand{\N}{\mathbb N}
\newcommand{\Z}{\mathbb Z}
\newcommand{\Q}{\mathbb Q}

% LTL
\newcommand{\true}{\mathbf{true}}
\newcommand{\false}{\mathbf{false}}
\newcommand{\X}{\mathbf X}
\newcommand{\U}{\mathbf U}
\newcommand{\Release}{\mathbf R}
\newcommand{\F}{\mathbf F}
\newcommand{\G}{\mathbf G}

% CTL
\newcommand{\EX}{\mathbf{EX}}
\newcommand{\AX}{\mathbf{AX}}
\newcommand{\EF}{\mathbf{EF}}
\newcommand{\AF}{\mathbf{AF}}
\newcommand{\EG}{\mathbf{EG}}
\newcommand{\AG}{\mathbf{AG}}
\newcommand{\EU}[2]{\mathbf E(#1 \U #2)}
\newcommand{\AU}[2]{\mathbf A(#1 \U #2)}
\newcommand{\ER}[2]{\mathbf E(#1 \Release #2)}
\newcommand{\AR}[2]{\mathbf A(#1 \Release #2)}

\newcommand{\NL}{\textsf{NL}}
\newcommand{\PTIME}{\textsf{PTIME}}
\newcommand{\PSPACE}{\textsf{PSPACE}}

\makeatletter

\def\dasharrowfill@#1#2#3#4{%
        $\m@th
        \thickmuskip0mu
        \medmuskip\thickmuskip
        \thinmuskip\thickmuskip
        \relax
        #4#1\mkern2mu
        \xleaders\hbox{$#4\mkern2mu#2\mkern2mu$}\hfill
        \mkern2mu
        #3$%
}

\def\dashleftarrowfill@{\dasharrowfill@\leftarrow\relbar\relbar}
\def\dashrightarrowfill@{\dasharrowfill@\relbar\relbar\rightarrow}
\def\dashleftrightarrowfill@{\dasharrowfill@\leftarrow\relbar\rightarrow}
\def\dashLeftarrowfill@{\dasharrowfill@\Leftarrow\Relbar\Relbar}
\def\dashRightarrowfill@{\dasharrowfill@\Relbar\Relbar\Rightarrow}
\def\dashLeftrightarrowfill@{\dasharrowfill@\Leftarrow\Relbar\Rightarrow}

\providecommand*\xdashleftarrow[2][]{%
  \ext@arrow 0055{\dashleftarrowfill@}{#1}{#2}}
\providecommand*\xdashrightarrow[2][]{%
  \ext@arrow 0055{\dashrightarrowfill@}{#1}{#2}}
\providecommand*\xdashleftrightarrow[2][]{%
  \ext@arrow 0055{\dashleftrightarrowfill@}{#1}{#2}}
\providecommand*\xdashLeftarrow[2][]{%
  \ext@arrow 0055{\dashLeftarrowfill@}{#1}{#2}}
\providecommand*\xdashRightarrow[2][]{%
  \ext@arrow 0055{\dashRightarrowfill@}{#1}{#2}}
\providecommand*\xdashLeftrightarrow[2][]{%
  \ext@arrow 0055{\dashLeftrightarrowfill@}{#1}{#2}}

\def\slashedarrowfill@#1#2#3#4#5{%
$\m@th\thickmuskip0mu\medmuskip\thickmuskip\thinmuskip\thickmuskip
  \relax#5#1\mkern-7mu%
  \cleaders\hbox{$#5\mkern-2mu#2\mkern-2mu$}\hfill
  \mathclap{#3}\mathclap{#2}%
  \cleaders\hbox{$#5\mkern-2mu#2\mkern-2mu$}\hfill
  \mkern-7mu#4$%
}

\def\rightslashedarrowfill@{%
  \slashedarrowfill@\relbar\relbar{\raisebox{1.2pt}{$\scriptscriptstyle\diagup$}}\rightarrow}
\newcommand\xslashedrightarrow[2][]{%
  \ext@arrow 0055{\rightslashedarrowfill@}{#1}{#2}}
\makeatother

\newtheorem{exercise}{Exercise}

\newenvironment{solution}{\begin{proof}[\protect{\textnormal{\emph{Solution.}}}]}{\end{proof}}

\let\oldqedhere\qedhere

\newif\ifstartedinmathmode
\renewcommand*{\qedhere}{%
  \relax\ifmmode\startedinmathmodetrue\else\startedinmathmodefalse\fi%
  \ifstartedinmathmode\tag*{\protect\oldqedhere}\else\ifinlist{\hfill\oldqedhere}{\oldqedhere}\fi%
}

\crefname{equation}{}{}
\crefname{exercise}{Exercise}{Exercises}


\title{Formal verification - Tutorial 02}

\begin{document}

\maketitle

\section*{Linear temporal logic}

Fix a set of propositional variables $P$.
%
Recall that linear temporal logic (LTL) formulas over $P$
are defined by the following (minimalistic) grammar:
%
\begin{align*}
	\varphi, \psi ::= \true \sep \false \sep p \sep \neg \varphi \sep \varphi \wedge \varphi \sep \X \varphi \sep \varphi \U \psi,
\end{align*}
%
where $p \in P$.
%
We define the derived operators
%
\begin{align*}
	\F \varphi &= \true \U \varphi \\
	\G \varphi &= \neg \F \neg \varphi.
\end{align*}

\begin{exercise}
	Write LTL formulas for the following properties:
	\begin{enumerate}
		\item $p$ holds infinitely often;
		\item whenever $p$ holds, then $q$ holds afterwards;
		\item $p$ holds \emph{precisely} at every even step;
	\end{enumerate}
\end{exercise}

\begin{solution}
	To express that $p$ holds infinitely often,
	we can write $\G \F p$.
	The second property is expressed by $\G (p \to \F q)$.
	The last property can be expressed in a variety of ways.
	For instance, we can write
	%
	\begin{align*}
		p \wedge \G ((p \implies \X \neg p) \wedge (\neg p \implies \X p)). \qedhere
	\end{align*}
\end{solution}

\begin{exercise}
	Consider the property
	``for every prefix of the computation,
	the number of times $p$ holds is at most the number of times $q$ holds''.
	Is this property expressible in LTL?
\end{exercise}

\begin{solution}
	No, the stated property is not expressible in LTL.
	The reason is that LTL formulas can only express $\omega$-regular properties,
	and the property in question is not $\omega$-regular.
\end{solution}

\begin{exercise}
	Consider the property
	``$p$ holds at all even time steps'' (but possibly also at odd steps).
	Is this property expressible in LTL?
\end{exercise}

\begin{solution}
	This property is not expressible in LTL,
	but in order to show it, we need to use a more advanced tool (next tutorial).
\end{solution}

\section*{From LTL to alternating Büchi automata}

We recall the translation from LTL to alternating Büchi automata.
%
To this end, we first convert an LTL formula $\varphi$ to negation normal form,
i.e., we push negations inside until they only occur in front of propositional variables.
%
To achieve this, we introduce the De Morgan dual of the until operator, which is called the \emph{release} operator:
%
\begin{align*}
	\varphi \Release \psi &= \neg (\neg \varphi \U \neg \psi).
\end{align*}
%
Note that $\G \varphi$ is equivalent to $\false \Release \varphi$.

We will use the following fixpoint expansions of the until and release operators:
%
\begin{align*}
	\varphi \U \psi &= \psi \vee (\varphi \wedge \X (\varphi \U \psi)), \\
	\varphi \Release \psi &= \psi \wedge (\varphi \vee \X (\varphi \Release \psi)).
\end{align*}

Given an LTL formula $\varphi$ in negation normal form,
we construct an alternating Büchi automaton
$A = \tuple{\Sigma, Q, q_0, F, \delta}$ as follows.
%
States of the automaton $Q$ are the subformulas of $\varphi$,
to which we add $\true$ and $\false$.
%
The initial state is $q_0 = \set \varphi$
and the set of accepting states $F$ contains all the release formulas and $\true$.
%
The transition function is defined by induction
on the structure of states $\psi \in Q$.
%
The cases for propositional variables and positive Boolean connectives are straightforward:
%
\begin{align*}
	\delta(\true, a) &= \true, \\
	\delta(\false, a) &= \false, \\
	\delta(p, a) &= \true \text{ if } p \in a, \text{ and } \false \text{ otherwise}, \\
	\delta(\neg p, a) &= \true \text{ if } p \not\in a, \text{ and } \false \text{ otherwise}, \\
	\delta(\psi \wedge \chi, a) &= \delta(\psi, a) \wedge \delta(\chi, a), \\
	\delta(\psi \vee \chi, a) &= \delta(\psi, a) \vee \delta(\chi, a).
\end{align*}
%
The most interesting cases are those for the temporal operators.
%
\begin{align*}
	\delta(\X \psi, a)
		&= \psi, \\
	\delta(\psi \U \chi, a)
		&= \delta(\chi, a) \vee (\delta(\psi, a) \wedge \psi \U \chi), \\
	\delta(\psi \Release \chi, a)
		&= \delta(\chi, a) \wedge (\delta(\psi, a) \vee \psi \Release \chi).
\end{align*}
%
For until and release we have applied the transition function $\delta(\cdot, a)$
to the fixpoint expansions given above.

\begin{exercise}
	Apply the translation above
	to construct an alternating Büchi automaton
	for the LTL formula $\G \F p$.
\end{exercise}

\begin{exercise}
	Using the automaton above, how can we construct an automaton for $\G \F p \wedge \G \F q$?
\end{exercise}

\end{document}
